\documentclass[11pt]{article}
\usepackage[T1]{fontenc}
\usepackage{lmodern,amsmath,amssymb,amsthm}
\usepackage[margin=1in]{geometry}
\usepackage{microtype}
\usepackage{xurl}
\usepackage[colorlinks=true,urlcolor=blue,linkcolor=blue,citecolor=blue]{hyperref}
\newtheorem{theorem}{Theorem}
\newtheorem{lemma}[theorem]{Lemma}
\newcommand{\C}{\mathbb C}
\newcommand{\R}{\mathbb R}
\newcommand{\N}{\mathbb N}
\newcommand{\J}{J}
\newcommand{\K}{K}
\title{Conjecture 00000003348:\\An expanding polynomial whose Julia set is a circle}
\author{C0ldSmi1e\thanks{AI-assisted contribution using OpenAI Codex. Separate AI review records accompany the submission; no external human review or maintainer acceptance is claimed.}}
\date{11 October 2026 (UTC)}
\begin{document}
\maketitle
\begin{abstract}
Under the standard complex-polynomial interpretation, the conjecture is false.
For $p(z)=z^2$, the filled Julia set defined by bounded forward orbits is the
closed unit disk, and its boundary is the unit circle. Every derivative of an
iterate expands every tangent vector on that circle by exactly $2^n$.
The circle is connected and has distinct points, so it is not homeomorphic to
Cantor space. This contradicts the conjecture's necessary Cantor-type clause;
no conclusion about the separate pressure--dimension formula is asserted.
The complete argument is formalized in Lean 4 with pinned Mathlib dependencies.
\end{abstract}

\section{The statement and the conventions being used}
The supplied English text is:
\begin{quote}
Definition: Expansive maps: uniformly expanding dynamics. Conjecture: The
embedding of expansive kernels: the Julia set of an expansive map is of Cantor
type, with the dimension of the Cantor set the root of the pressure equation.
(expansive Cantor pressure dimension root)
\end{quote}
The Chinese text likewise specifies uniform expansion and asserts both
Cantor type and a dimension--pressure conclusion. Both exact texts are
preserved in \texttt{original.md}. Neither text specifies an ambient class of
maps, an expansion locus or metric, a Julia construction, a dimension
convention, or a pressure potential. We disclose the following standard
complex-dynamical reading, rather than claiming that the text determines a
unique formal proposition.

We work with polynomials $p\in\C[X]$ of degree at least two, acting on $\C$
with its usual topology and Euclidean norm. Write $f=p.\mathrm{eval}$ and
$f^{\circ 0}=\mathrm{id}$. Define
\[
 \K(p)=\{z\in\C:\exists R\in\R\ \forall n\in\N,
                   |f^{\circ n}(z)|\le R\},\qquad
 \J(p)=\partial\K(p).
\]
This is the standard filled-Julia and boundary-Julia convention for
polynomials of degree at least two; see Cheraghi~\cite{cheraghi}, Lecture~4,
p.~22, Lemma~4.1. Lean uses this boundary definition directly. It does not
formalize normal families or newly prove the equivalence with that alternative
Julia definition.

Uniform expansion on a compact invariant set $S$ means that there are
$C>0$ and $L>1$ with
\[
 \|D(f^{\circ n})_z(v)\|\ge C L^n\|v\|
 \quad(n\ge0,\ z\in S,\ v\in T_z\C).
\]
This is the usual tangent-vector definition; see Milnor~\cite{milnor},
\S14, p.~14-1. Our Lean predicate additionally requires differentiability of
every iterate on $\C$ and expresses the bound by the complex derivative norm.
For the example we also prove the corresponding real tangent-plane equality,
so this scalar presentation is connected to the actual derivative acting on
vectors. Expansion is asserted on $\J(p)$, as is conventional for expanding
polynomial dynamics. It is not asserted on the entire plane or sphere.

Cantor type means homeomorphism of the entire Julia subspace with the standard
Cantor space $\{0,1\}^{\N}$, with its product topology. It does not mean merely
containing an embedded Cantor subset. The introductory phrase about
``embedding of expansive kernels'' supplies no separately defined object or
condition; we infer no hidden disconnectedness or inverse-branch separation
assumption from it. A different intended ambient category or an additional
unstated hypothesis would require a different formal statement.

\section{The counterexample}
\begin{theorem}
For $p(X)=X^2$, the Julia set $\J(p)$ is compact and forward invariant,
$p$ is uniformly expanding on it, and $\J(p)$ is not of Cantor type.
\end{theorem}
\begin{proof}
The polynomial has degree two and evaluation $f(z)=z^2$. Induction gives
\[
 f^{\circ n}(z)=z^{2^n},\qquad |f^{\circ n}(z)|=|z|^{2^n}
 \quad(n\in\N).
\]
If $|z|\le1$, all these norms are at most one. If $|z|>1$, then for any
proposed bound $R$, unboundedness of ordinary powers gives $n$ with
$R<|z|^n$. Since $n\le2^n$ and the base exceeds one,
\[
 R<|z|^n\le |z|^{2^n}=|f^{\circ n}(z)|,
\]
contradicting that bound. Thus
\[
 \K(p)=\{z:|z|\le1\},\qquad
 \J(p)=\partial\K(p)=\{z:|z|=1\}.
\]
The last equality is the actual topological boundary of a closed ball of
positive radius in the complex plane. The resulting sphere is compact.
It is forward invariant because $|z^2|=|z|^2=1$ on it.

Every iterate is complex differentiable and
\[
 (f^{\circ n})'(z)=2^n z^{2^n-1}.
\]
For $|z|=1$, its norm is $2^n$. The complex derivative acts by scalar
multiplication by $(f^{\circ n})'(z)=2^n z^{2^n-1}$, hence for every complex vector $v$,
\[
 \|D_{\C}(f^{\circ n})_z(v)\|=2^n\|v\|.
\]
Restriction of scalars gives the same equality for the real derivative on
the underlying real plane. This includes $n=0$ and proves uniform expansion
with the common constants $C=1$ and $L=2$.

The unit sphere in a real normed space of dimension greater than one is
connected; the complex plane has real dimension two. Thus $\J(p)$ is
connected. Both $1$ and $-1$ belong to it and are distinct. If $\J(p)$ were
totally disconnected, its connected subset equal to itself would be a
subsingleton, a contradiction. Finally, Cantor space is totally disconnected,
and a homeomorphism transports that property. Therefore no homeomorphism
$\J(p)\cong\{0,1\}^{\N}$ exists.
\end{proof}

For completeness, the formalization also proves a directly metric expansion
estimate. If $|z|=1$ and $|z-w|\le\tfrac12$, then the triangle inequality gives
\[
 2=|2z|\le |z+w|+|z-w|,\qquad |z+w|\ge\tfrac32,
\]
and consequently
\[
 |f(z)-f(w)|=|z-w|\,|z+w|\ge\tfrac32|z-w|.
\]
The constants are uniform; $w$ may be any complex point in this neighbourhood.
This is a local statement: $f(1)=f(-1)$ rules out a global distance-expansion
inequality on every pair of Julia points.

\section{Why this disproves the stated conjunction}
The example satisfies the degree and uniform-expansion hypotheses in the
declared standard setting but fails the Cantor-type conclusion. Hence
\[
 \neg\,\forall p,\quad
 [\deg p\ge2\ \wedge\ p\text{ uniformly expanding on }\J(p)]
 \ \Longrightarrow\ \J(p)\text{ of Cantor type}.
\]
Adding any second conclusion to the right-hand side cannot make this
universal assertion true. The Lean corollary makes that logical point for
every predicate $A:\C[X]\to\mathrm{Prop}$: the conjunction
``$\J(p)$ is of Cantor type and $A(p)$'' still fails for the example.
This arbitrary predicate is not a substitute definition of pressure or
dimension. Neither quantity is computed, and the pressure--dimension identity
itself is not disproved. One false necessary conjunct suffices.

The origin provides a useful scope check: it belongs to $\K(p)$ but not
$\J(p)$, and $f'(0)=0$. Thus the example would not satisfy a different
condition requiring expansion on all of $\C$. The submitted interpretation
does not claim that it does.

\section{Formal correspondence and verification}
The project uses Lean 4.19.0 and Mathlib v4.19.0, with the exact dependency
revisions recorded in \texttt{lake-manifest.json}. All mathematical declarations
are in the namespace \path{ExpandingJuliaCounterexample}. The main file is
\texttt{Counterexample.lean}; \texttt{Checks.lean} contains boundary checks and
15 named axiom inspections.

\begin{itemize}
\item \path{filledJulia_square} and \path{julia_square} derive the disk
and circle from genuine polynomial iteration, orbit boundedness, and frontier.
\item \path{square_iterate_hasDerivAt} proves the actual derivative;
\path{norm_real_fderiv_square_iterate} supplies the real tangent equality;
\path{square_uniformlyExpanding} packages compactness, invariance,
differentiability, and uniform constants.
\item \path{cantorType_totallyDisconnected} connects the genuine
homeomorphism definition with total disconnectedness.
\path{julia_square_not_cantorType} proves its failure for the actual Julia set.
\item \texttt{counterexample} is the existential capstone. The final
conjunction corollary is
\path{not_every_expanding_polynomial_satisfies_conjunction}.
\end{itemize}

The author build passed with warnings treated as errors. The named axiom
inspections use only the standard foundational axioms
\texttt{propext}, \texttt{Classical.choice}, and \texttt{Quot.sound}.
There are no admitted proofs, added axioms, or uses of \texttt{native\_decide}
in the submitted mathematical source. Boundary checks cover the origin,
the zeroth iterate, and noninjectivity at $1,-1$.
No numerical experiment or auxiliary computational identity is needed for
the mathematics. The executable verification tools and their regression
controls are documented separately in \texttt{VERIFICATION.md}, along with
fresh replay, dependency inspection, exact hashes, and independent AI reviews.
Source compilation provides Lean kernel checking; a dependency inspector
is not an independent kernel implementation. These are local validation
records, not hosted CI success or maintainer acceptance. The example is
classical; no claim of a new mathematical discovery is made.

\begin{thebibliography}{9}
\bibitem{cheraghi}
Davoud Cheraghi, \emph{One dimensional real and complex dynamics}, lecture
notes, 8 December 2014, Lecture 4, p.~22, Lemma~4.1.
\url{https://www.ma.imperial.ac.uk/~dcheragh/Research/CourseNotes.pdf}.
\bibitem{milnor}
John Milnor, \emph{Dynamics in One Complex Variable}, course copy,
\S14, p.~14-1, definition of expansion on a compact invariant set.
\url{https://abel.math.harvard.edu/archive/118r_spring_05/docs/milnor.pdf}.
\end{thebibliography}
\end{document}
